\chapter{System Design and Methodology}
\label{ch:design}
\vspace{2em}

% TODO: Introduce the chapter -- summarise the design goals of the
% AI-Powered Omniscient Teacher Assistant and how they address the problem
% statement (administrative task-drift, feedback gap, clarity deficit).

\section{Design Goals and Requirements}
\label{design:sec:requirements}

% TODO: Functional and non-functional requirements derived from the problem
% statement and stakeholder needs (faculty, students, administrators).

\section{Deployment Model Analysis}
\label{design:sec:deployment}

% TODO: Frame the analysis -- before committing to an architecture, the system's
% deployment model must be chosen. Introduce the three candidates (fully cloud,
% fully on-premises, hybrid) and state that they are assessed against the
% requirements from \autoref{design:sec:requirements}.

\autoref{design:tab:deployment} compares the three candidate deployment models
across the dimensions most relevant to a virtual teaching assistant.

\begin{table}[htb]
  \centering
  \caption{Comparison of deployment models for the virtual teaching assistant.}
  \label{design:tab:deployment}
  \footnotesize
  \renewcommand{\arraystretch}{1.3}
  \begin{tabular}{L{0.13\textwidth} L{0.18\textwidth} L{0.18\textwidth} L{0.18\textwidth} L{0.185\textwidth}}
    \hline
    \textbf{Dimension} & \textbf{Cloud} & \textbf{On-Premises} & \textbf{Hybrid} & \textbf{Relevance to This System} \\
    \hline
    Scalability &
      Auto-scaling, elastic resources &
      Fixed hardware capacity &
      Burst to cloud for peak load &
      Needed when many students interact simultaneously \\
    Initial cost (CapEx) &
      Low upfront cost &
      High hardware investment &
      Moderate; reuses existing hardware &
      Important for student projects with limited budget \\
    Operational cost (OpEx) &
      Pay-as-you-go pricing &
      Maintenance, electricity, upgrades &
      Both cost models apply &
      Cloud easier to forecast for experimentation \\
    AI/ML infrastructure &
      Managed GPUs, LLM APIs, training services &
      GPUs and frameworks configured manually &
      Local inference; managed services for heavy workloads &
      Faster iteration for adaptive learning models \\
    Maintenance and updates &
      Managed by provider &
      Self-managed patches and upgrades &
      Split responsibility &
      Frees effort for algorithm design over operations \\
    Security controls &
      Built-in IAM, encryption, monitoring &
      Full control, higher responsibility &
      Must secure both environments and the link &
      Student learning data requires strong protection \\
    Data privacy and control &
      Shared responsibility model &
      Full ownership of data &
      Sensitive data on-premises; compute in cloud &
      On-premises stronger for sensitive datasets \\
    Reliability and availability &
      Multi-region redundancy &
      Depends on local setup &
      Cloud fallback for local outages &
      Important for continuous student access \\
    Latency &
      Depends on region selection &
      Very low on local network &
      Low for local paths; variable for cloud paths &
      Useful for real-time AI responses \\
    Development speed &
      Fast prototyping with managed services &
      Longer setup time &
      Added integration complexity &
      Accelerates the project timeline \\
    Vendor lock-in risk &
      Higher (managed services) &
      Lower &
      Lower; workloads remain portable &
      Key architectural consideration \\
    \hline
  \end{tabular}
\end{table}

The system adopts a \textbf{hybrid} deployment. Response generation is served
by open-weight models (OpenAI's gpt-oss family) on a hosted inference
provider~\cite{groq26:Models}. Embedding and retrieval remain local, using the
BGE-M3 embedder~\cite{arxiv24:Chen} on the operator's own hardware. A small local
model (Llama~3.2, 3B parameters) serves as a generation fallback. Student records,
course material and the vector index stay in the system's own database throughout.
The decision rests on measurements taken on the deployed system rather than on
the generic properties in \autoref{design:tab:deployment}; the full comparison is
reported in \autoref{eval:sec:results:providers}.

\paragraph{Why hosted generation.} On the system's flagship multi-part
interaction (look up the student's progress, then set a practice question),
hosted generation reduced median turn latency from 11.2\,s with the most capable
local model (Qwen3 14B) to 1.7\,s. Grounding fidelity did not suffer: the
assistant reported the student's actual figures in 10 of 10 turns in both cases.
Hosted generation also produced correct quiz answer keys more than twice as often
as the local quiz model. Its marginal cost is small: US\$0.0003--0.0005 per agent
turn, an estimated US\$0.10--0.27 per student per semester.

\paragraph{Why not fully on-premises.} On the 32\,GB host, the only local model
that grounded reliably (Qwen3 14B, 10.4\,GB resident) competed with the retriever
for memory. When loaded alongside the embedding model it evicted it, so every
subsequent retrieval paid a cold reload. The smaller local model that fits beside
the retriever grounded less reliably (8 of 10 turns) and produced correct quiz
answer keys in only 3 of 10 generated questions. A fully local deployment would
therefore have forced a choice between latency, retrieval performance and answer
quality.

\paragraph{Why not fully cloud.} The chosen inference provider offers no embedding
models, so retrieval cannot move without also changing the embedder, and with it
re-indexing the corpus. More importantly, a local fallback turns provider-side
failures into a drop in quality rather than an outage. Such failures include
outages, rate limiting and the retirement of hosted model identifiers, which
disabled quiz generation in production once during development. The fallback
only engages before any output has been streamed, so a student never receives an
answer spliced together from two models.

\paragraph{Trade-offs accepted.} Two costs of this choice are acknowledged.
First, the \emph{data-privacy} advantage of on-premises deployment is only
partly retained. Stored student records never leave the system's own database,
but each generation request carries the prompt for that turn to the inference
provider. That prompt includes the retrieved course excerpts and, when the
assistant consults the student's progress, figures such as quiz scores. Second,
the fallback path depends on a single local host being reachable, so the
redundancy it provides is limited to that host's availability.

% TODO: tie these deciding factors to the specific requirements in
% \autoref{design:sec:requirements} once that section is written (latency,
% grounding fidelity, budget, and data handling map directly onto them).

\section{System Architecture}
\label{design:sec:architecture}

% TODO: High-level architecture diagram and component overview
% (e.g. retrieval-augmented generation pipeline, knowledge base, analytics).

\section{Methodology}
\label{design:sec:methodology}

% TODO: Research methodology, data sources, and the approach taken to
% evaluate the system.

\subsection{Baseline Token Audit}
\label{design:subsec:tokenaudit}

The literature reviewed in \autoref{review:subsec:tokens} establishes in
general that a RAG pipeline's input token count is dominated by retrieved
context~\cite{femnlp23:Liu}; this section establishes it for \emph{this}
system, in its own data, because that measured share --- not the citation ---
is the design motivation for the compression stage. The audit decomposes the
uncompressed pipeline's prompt into its four components:

\begin{equation}
  \label{design:eq:tokenaudit}
  T_{\mathrm{input}}
    \;=\;
  \underbrace{k \cdot \bar{t}_{\mathrm{chunk}}}_{\text{retrieved context}}
  \;+\; t_{\mathrm{system}}
  \;+\; t_{\mathrm{history}}
  \;+\; t_{\mathrm{query}},
\end{equation}

where $k$ is the number of retrieved chunks per query,
$\bar{t}_{\mathrm{chunk}}$ the mean tokens per chunk under the deployed
Canvas-material chunking configuration, and $t_{\mathrm{system}}$,
$t_{\mathrm{history}}$, and $t_{\mathrm{query}}$ the system prompt,
conversation history, and student question respectively. The quantities of
interest are the baseline input tokens per query and the share attributable
to the retrieved-context term. This baseline anchors the thesis's coherence
loop: the problem statement quantifies cost in tokens per query
(\autoref{intro:sec:problem}), this audit measures the system's baseline in
the same units, and the evaluation reports compression's effect in the same
units again (\autoref{eval:sec:metrics}).

% TODO: run the audit on the actual pipeline (measure tokens per chunk and k
% from the deployed chunking config; take system prompt / history / query
% tokens from production logs) and report the numbers and the
% retrieved-context share here.

\subsection{Selection of the Compression Method}
\label{design:subsec:compression-choice}

Among the compression methods surveyed in \autoref{review:subsec:compression},
this thesis adopts the extractive, sentence-level variant of
\textsc{Recomp}~\cite{iclr24:Xu}, for four reasons.

First, fit to the data: course materials are sentence-dense expository text,
and extractive sentence selection preserves the instructor's exact wording
--- important for a study assistant that must not paraphrase away technical
definitions --- whereas abstractive summarisation risks hallucinated
rewording of precisely the content students ask about.

Second, fit to the deployment target: the extractive compressor requires only
dual-encoder sentence scoring, which is feasible on consumer-grade local
hardware alongside the generation model itself. The alternatives embed a
language model in the compression loop --- the LLMLingua family uses an
auxiliary LM for perplexity scoring~\cite{emnlp23:Jiang, acl24:Jiang}, and
ECoRAG relies on LLM-mined evidentiality labels~\cite{facl25:Jeong} --- which
the memory- and compute-constrained host of this system cannot afford.

Third, fit to the pipeline: compression is architecturally a drop-in stage
between retriever and generator, leaving both untouched --- matching the
modular design of \autoref{design:sec:architecture} and enabling the
feature-flagged A/B comparison used in \autoref{ch:evaluation}.

Fourth, evidence of headroom: across the surveyed methods, QA-style workloads
consistently admit 50--95\% token reduction at negligible quality
cost~\cite{femnlp23:Liu, iclr24:Xu, emnlp23:Jiang, acl24:Jiang, emnlp23:Li},
and the material-grounded question answering of a VTA is exactly that
workload shape.

Finally, the choice carries the novelty position stated in
\autoref{review:subsec:compression}: no published work evaluates
\textsc{Recomp} in a deployed educational VTA, so evaluating it here is
itself a contribution rather than a re-implementation. The realisation of
the stage is described in \autoref{impl:sec:compression}.

\section{Summary}
\label{design:sec:summary}

% TODO: Recap the design decisions that the following chapter implements.
